\section{方法}
\subsection{框架概述}
ReCA-SAM3 由 SAM 3 特征提取、RemoteCLIP 语义引导、双流特征聚合和语义—细节融合四部分组成，如图~\ref{fig:overview} 所示。SAM 3 为各类别提示提取条件特征，RemoteCLIP 提供局部图文相似度先验。聚合模块在低分辨率网格上交替建模提示间关系与提示内空间联系，再通过语义—细节融合（Semantic--Detail Fusion，SDF）将更新特征接入冻结的解码器。

\begin{figure*}[!t]
\centering
\fbox{\begin{minipage}{0.93\textwidth}
\centering\small
Image and category prompts $\longrightarrow$ SAM 3 prompt-conditioned features\\
RemoteCLIP image--text similarities $\longrightarrow$ Semantic guidance\\
Dual-stream cross-class and spatial aggregation $\longrightarrow$ Semantic--Detail Fusion\\
Frozen SAM 3 pixel decoder $\longrightarrow$ Per-class semantic predictions
\end{minipage}}
\caption{Overview of ReCA-SAM3. RemoteCLIP guides dual-stream feature aggregation before the frozen SAM 3 decoder. Semantic--Detail Fusion combines aggregated features with the original semantic and spatial representations.}
\label{fig:overview}
\end{figure*}

\subsection{任务定义与提示映射}
给定遥感图像和候选类别名称，模型输出每个像素的语义标签。包含多个同义名称的类别被展开为独立文本提示，各提示共享原始类别监督。聚合在同一图像的全部展开提示上执行，随后分块解码。推理时，对属于同一类别的提示概率逐像素取最大值，得到类别响应，再进行类别间决策。

类间注意力沿展开提示轴交互，因此同时包含不同类别和同类别别名之间的信息交换。类内注意力仅在单个提示的空间网格上计算。两种交互分别对应提示语义联系与空间结构，不混合不同图像。

\subsection{RemoteCLIP 语义引导}
RemoteCLIP 为图像提取局部视觉特征，并使用遥感文本模板编码类别名称~\cite{liuRemoteCLIPVisionLanguage2024}。归一化的局部视觉特征与模板文本特征计算余弦相似度，形成逐提示分数图。分数图经通道投影后，与局部视觉特征分别归一化、拼接并融合；随后通过空间卷积生成初始语义状态。

SAM 3 的提示条件特征被下采样为初始特征状态。两条状态流均位于 $36\times36$ 网格，具有 256 个通道。RemoteCLIP 的中间层视觉特征为语义流的类内注意力提供指导，SAM 3 特征金字塔网络（Feature Pyramid Network，FPN）的空间特征为特征流提供指导。

\subsection{双流特征聚合}
双流聚合采用类间块与类内块交替的结构。类间块在同一位置沿提示轴计算注意力，依次更新语义状态与特征状态。每次更新的查询和键由两条状态共同构造，值来自被更新的状态，并通过残差连接保留原有表示。

类内块在单个提示内依次执行普通窗口与移位窗口注意力，先更新语义状态，再更新特征状态。对应的视觉指导特征参与查询和键的构造。四个聚合块按类间、类内、类间、类内排列，块间不共享参数；各块末尾不使用额外前馈网络。聚合结束后，对特征状态进行通道归一化，并送入 SDF。

\subsection{语义—细节融合与解码}
SDF 将聚合特征上采样至 $72\times72$，分别与原始提示条件特征及 FPN 空间特征构成语义支路和细节支路。各输入经通道投影后相加，每条支路使用自身的共享卷积块处理 $0^\circ$、$90^\circ$、$180^\circ$ 和 $270^\circ$ 四个旋转视图。卷积输出逆旋转至原坐标后取平均，以残差形式更新支路表示。

两条支路融合后生成解码特征，替换 SAM 3 Pixel Decoder 的最后一级输入，其余高分辨率 FPN 特征保持原解码路径。SAM 3 参数冻结，学生分支中的梯度仍可穿过解码器回传至 SDF 和聚合模块。RemoteCLIP 的图像与文本分支仅微调注意力 Q/V 投影及位置参数。

\subsection{训练目标与推理}
训练采用二元交叉熵分割监督，并结合原始 SAM 3 响应蒸馏与存在类别的忽略区域抑制。参考分支使用原始提示条件特征和相同的冻结解码器，在无梯度条件下产生软掩码。分割损失作用于有效标签像素；蒸馏结合有效区域与指定边界邻域；抑制项约束存在类别在忽略区域的响应。同义提示按其原始类别组织监督。

推理先将提示响应合并到语义类别，再执行参考概率门与预测概率门。参考门依据原始 SAM 3 响应过滤类别分数，预测门依据学生响应过滤低置信度区域。类别决策与背景回填使用同一标签映射。测试时增强关闭。

